\documentclass[11pt,a4paper]{article}
\usepackage[margin=25mm,headheight=14pt]{geometry}
\usepackage{fontspec}
\setmainfont{lmroman10-regular.otf}[BoldFont=lmroman10-bold.otf,ItalicFont=lmroman10-italic.otf,BoldItalicFont=lmroman10-bolditalic.otf]
\setsansfont{lmsans10-regular.otf}
\setmonofont{lmmono10-regular.otf}[Scale=MatchLowercase]
\newfontfamily\authorsymbol{Noto Sans Symbols}[Scale=MatchUppercase]
\usepackage{amsmath,amssymb,amsthm,mathtools,booktabs,array}
\usepackage{microtype}
\usepackage[hidelinks,unicode]{hyperref}
\usepackage{fancyhdr}
\pagestyle{fancy}\fancyhf{}
\fancyhead[L]{\small\itshape Sum-and-carry functions for cascaded addition}
\fancyhead[R]{\small\thepage}\renewcommand{\headrulewidth}{0pt}
\setlength{\parindent}{1.2em}\setlength{\parskip}{0pt}
\setlength{\emergencystretch}{2em}
\newtheorem{theorem}{Theorem}
\newtheorem{proposition}[theorem]{Proposition}
\newtheorem{lemma}[theorem]{Lemma}
\newtheorem{corollary}[theorem]{Corollary}
\theoremstyle{definition}\newtheorem{definition}[theorem]{Definition}
\newcommand{\F}{\mathbb F_2}
\newcommand{\andb}{\mathbin{\&}}
\newcommand{\res}[1]{[#1]_M}
\DeclareMathOperator{\wt}{wt}
\DeclareMathOperator{\rank}{rank}
\hypersetup{pdftitle={Exact Counts of Sum-and-Carry Functions for Cascaded Addition},pdfauthor={lmtlssss🜐},pdfsubject={Mathematics: enumerative combinatorics and arithmetic functions},pdfkeywords={carry functions, cascaded addition, exact enumeration, minimal encoding, partial permutations}}
\title{\LARGE Exact Counts of Sum-and-Carry Functions\\for Cascaded Addition}
\author{lmtlssss{\authorsymbol 🜐}}
\date{27 October 2026\\[3pt]{\small\href{https://doi.org/10.5281/zenodo.23700193}{doi:10.5281/zenodo.23700193}}}
\newcommand{\Q}{\mathcal Q}
\newcommand{\supp}{\mathcal R}
\newcommand{\ind}{\mathbf 1}
\newcommand{\abits}{\boldsymbol a}
\newcommand{\xvec}{\boldsymbol x}
\begin{document}
\maketitle\thispagestyle{plain}
\begin{abstract}
We count the complete functions defined by successive radix-$B$ additions
with independently variable incoming carry bits. Each function retains the
final residue and every outgoing carry. For $k$ additions, let
$A_0=1$, $A_j=jA_{j-1}+1$, and $t=\min\{k,B-1\}$.
The exact number of functions is
$B^{k-t}((B-1)A_t+1)$.
The proof classifies each carry function by its possible final residues:
a full initial interval or a single terminal residue.
This classification gives a dense encoder and decoder, and hence the exact
minimum fixed-length label. Two additions give $5B-4$ functions; three give
$16B-15$ when $B\geq4$.
For three 32-bit additions, 36 bits preserve the complete function of four
32-bit operands, including all eight incoming-carry assignments.
\end{abstract}
\noindent\textit{Keywords.} Arithmetic carries; function classification;
exact enumeration; minimal codes; partial permutations.

\section{The arithmetic function}
Fix integers $B\geq2$ and $k\geq0$. Choose operands
$\xvec=(x_0,\ldots,x_k)\in\{0,\ldots,B-1\}^{k+1}$.
The operands are parameters. The variable inputs are independent bits
$\abits=(a_1,\ldots,a_k)\in\{0,1\}^k$.
Starting from $r_0=x_0$, define, for $1\leq i\leq k$,
\begin{equation}\label{eq:operation}
 u_i=r_{i-1}+x_i+a_i,\qquad
 c_i=\left\lfloor\frac{u_i}{B}\right\rfloor,\qquad
 r_i=u_i-Bc_i.
\end{equation}
Every $c_i$ is a bit. We retain the function
\begin{equation}\label{eq:function}
 Q_{\xvec}:\abits\longmapsto
      (r_k(\abits),c_1(\abits),\ldots,c_k(\abits)).
\end{equation}
Two operand tuples are equivalent when these functions agree on all
$2^k$ inputs. Write $\Q_k(B)$ for the set of distinct functions, and
$N_k(B)=|\Q_k(B)|$. For $k=0$, the function is the constant residue $x_0$,
so $N_0(B)=B$.

The distinction between parameters and inputs is essential. A label for
$Q_{\xvec}$ identifies an arithmetic function. It can represent several
operand tuples, all with the same responses to the incoming bits.
The specified outputs include each stage's outgoing carry, rather than only
their sum or the final modular result.

\section{Exact count}
Let
\begin{equation}\label{eq:arrangements}
 A_0=1,\qquad A_j=jA_{j-1}+1
       =j!\sum_{\ell=0}^{j}\frac1{\ell!}.
\end{equation}
These are the classical numbers of ordered subsets of a $j$-element set,
recorded as A000522~\cite{oeis}. Their initial values are
$1,2,5,16,65,326,1957$.

\begin{theorem}\label{thm:main}
For every $B\geq2$ and $k\geq0$, put $t=\min\{k,B-1\}$. Then
\begin{equation}\label{eq:count}
 \boxed{N_k(B)=B^{k-t}\bigl((B-1)A_t+1\bigr).}
\end{equation}
In particular, for $0\leq k<B$,
\begin{equation}\label{eq:short}
 N_k(B)=(B-1)A_k+1,
\end{equation}
and $N_k(B)=B\,N_{k-1}(B)$ for $k\geq B$.
\end{theorem}

The first cases are $2B-1$, $5B-4$, $16B-15$, and $65B-64$,
within the respective ranges $k<B$. The assertion counts functions,
including those attained by many different operand tuples.

\subsection{Cumulative carry profiles}
Set
\[
 S_i=x_0+\cdots+x_i,\qquad T_i=a_1+\cdots+a_i,
 \qquad P_i(t)=\left\lfloor\frac{S_i+t}{B}\right\rfloor
 \quad(0\leq t\leq i),
\]
with $P_0(0)=0$.
Repeated substitution in \eqref{eq:operation} gives
\begin{equation}\label{eq:prefix}
 r_i=S_i+T_i-B\sum_{j=1}^{i}c_j,
 \qquad \sum_{j=1}^{i}c_j=P_i(T_i).
\end{equation}
Hence
\begin{equation}\label{eq:reconstruct}
 c_i=P_i(T_i)-P_{i-1}(T_{i-1}),\qquad
 r_k=[v+T_k]_B,\qquad v=[S_k]_B,
\end{equation}
where $[z]_B$ is the least nonnegative residue modulo $B$.

\begin{lemma}\label{lem:profiles}
The profiles $(P_1,\ldots,P_k)$ determine the complete carry-output
function. Conversely, that function determines every profile.
Together with $v$, the profiles determine $Q_{\xvec}$ uniquely.
\end{lemma}
\begin{proof}
Equation~\eqref{eq:reconstruct} gives one direction.
For the other, choose a prefix of $i$ incoming bits with exactly $t$ ones.
The sum of its first $i$ carry outputs equals $P_i(t)$ by
\eqref{eq:prefix}. Every $t\in\{0,\ldots,i\}$ occurs.
The final residue at the all-zero input gives $v$.
\end{proof}

Write $S_i=Bq_i+v_i$. Then
\begin{equation}\label{eq:profile-form}
 P_i(t)=q_i+\left\lfloor\frac{v_i+t}{B}\right\rfloor.
\end{equation}
For $i<B$, this profile is constant exactly when
$0\leq v_i\leq B-i-1$. Otherwise its first increase occurs at
$t=B-v_i\in\{1,\ldots,i\}$ and identifies $v_i$.

\subsection{Residue supports}
For a carry-output function $\phi$ of $i$ incoming bits, let
$\supp(\phi)$ be the set of residues $[S_i]_B$ attained by operand tuples
with that same carry function.

\begin{lemma}\label{lem:support}
For $0\leq i<B$, every carry function has one of the following supports:
\begin{enumerate}
\item the interval $\{0,\ldots,B-i-1\}$;
\item a singleton $\{v\}$ with $B-i\leq v\leq B-1$.
\end{enumerate}
There are $A_i$ functions of the first kind. For each terminal residue,
there are $E_i=A_i-A_{i-1}$ functions of the second kind when $i\geq1$.
\end{lemma}
\begin{proof}
At $i=0$, the empty carry function has every residue as its support.
Write $A_0=1$ and $E_0=0$ for this initial case.
Assume the classification at $i-1$, and append an operand $h$.
Choose an old residue $r$ and a new residue $v$.
Exactly one $h\in\{0,\ldots,B-1\}$ satisfies $[r+h]_B=v$.
Its zero-input carry is
\begin{equation}\label{eq:e}
 e=\ind_{\{r>v\}}.
\end{equation}
If $q=P_{i-1}(0)$, the appended cumulative profile is
\begin{equation}\label{eq:extension}
 P_i(t)=q+e+\left\lfloor\frac{v+t}{B}\right\rfloor.
\end{equation}
Thus the old carry function, $v$, and $e$ determine the extension.
Different old functions remain distinct after extension; for a fixed old
function and $v$, the value at the zero input distinguishes the two values
of $e$.

Suppose $0\leq v\leq B-i-1$. Each old interval support
$\{0,\ldots,B-i\}$ permits both $e=0$ and $e=1$.
Every old singleton support has $r\geq B-i+1>v$, so it permits only $e=1$.
Equation~\eqref{eq:extension} is constant in $t$ and independent of the
chosen $v$ throughout this interval. Every resulting function therefore
attains the entire interval. The number of these carry functions is
\begin{equation}\label{eq:rec-a}
 A_i=2A_{i-1}+(i-1)E_{i-1}.
\end{equation}

Now let $B-i\leq v\leq B-1$. An old interval support permits only $e=0$.
Each old singleton permits exactly the value in \eqref{eq:e}.
The new profile has its first increase at $B-v$, so it identifies $v$
and has no other residue in its support. For each such $v$ the count is
\begin{equation}\label{eq:rec-e}
 E_i=A_{i-1}+(i-1)E_{i-1}.
\end{equation}
These constructions realize every listed extension with genuine operands.
Conversely, every operand extension uses one of the listed choices.
They are distinct by the preceding recovery argument.

Subtracting \eqref{eq:rec-e} from \eqref{eq:rec-a} gives
$E_i=A_i-A_{i-1}$. The initial values are $A_1=2$ and $E_1=1$.
Induction now yields $A_i=iA_{i-1}+1$:
for $i\geq2$, substitute $E_{i-1}=A_{i-1}-A_{i-2}$ in
\eqref{eq:rec-a}, and use
$A_{i-1}=(i-1)A_{i-2}+1$.
This is precisely \eqref{eq:arrangements}.
\end{proof}

\begin{proof}[Proof of Theorem~\ref{thm:main}]
For $1\leq k<B$, Lemma~\ref{lem:support} gives
\[
 N_k(B)=(B-k)A_k+kE_k
       =BA_k-kA_{k-1}=(B-1)A_k+1.
\]
The case $k=0$ was established above.

If $i\geq B-1$, a carry function determines its final residue.
Indeed, the values of $P_i(t)-P_i(0)$ for $0\leq t\leq B-1$
identify $v_i$: a first increase at $t$ gives $v_i=B-t$,
and no increase gives $v_i=0$.
Now take $k\geq B$. Every old function at length $k-1$ has a uniquely
determined residue $r$. Its $B$ possible appended operands give $B$
different final residues. Two different old functions remain distinct,
since the new output retains all old carries.
Therefore $N_k(B)=B N_{k-1}(B)$.
Iterating from $k=B-1$ proves \eqref{eq:count}.
\end{proof}

\section{Dense labels and the exact bit minimum}
\begin{theorem}\label{thm:code}
There is an explicit bijection between $\Q_k(B)$ and
$\{0,\ldots,N_k(B)-1\}$. Consequently the least number of bits in a
fixed-length exact function label is
\begin{equation}\label{eq:bits}
 b_k(B)=\left\lceil\log_2 N_k(B)\right\rceil.
\end{equation}
A decoder returns a representative operand tuple with the labelled
function, preserving every output in \eqref{eq:function}.
\end{theorem}
\begin{proof}
We specify a recursive ordering. At length zero, the label is the residue.
For $1\leq i<B$, order the functions first by their final residue $v$.
The offset of its group is
\begin{equation}\label{eq:offset}
 O_i(v)=
 \begin{cases}
 vA_i,&v<B-i,\\
 (B-i)A_i+(v-B+i)E_i,&v\geq B-i.
 \end{cases}
\end{equation}
Within a residue group, list old interval-support functions first.
If $v<B-i$, an old local index $p$ contributes the two indices
$2p$ and $2p+1$, corresponding to $e=0$ and $e=1$.
If $v\geq B-i$, it contributes only index $p$.
Next list the old singleton supports in increasing residue order, and
within each use the previous local ordering. Each supplies one extension.
Lemma~\ref{lem:support} proves that these lists have respectively
$A_i$ or $E_i$ entries, with unique coverage.

The ordering is constructive in both directions.
For a chosen interval-support predecessor, use representative residue
$r=0$ when $e=0$ and $r=v+1$ when $e=1$.
Both choices lie in the required old support.
For a singleton predecessor, use its determined residue.
Decode that predecessor recursively, then append $h=[v-r]_B$.
This realizes the prescribed extension.

For $i\geq B$, an old label $z$ and appended operand $h$ have new label
$Bz+h$. Quotient and remainder recover them. The old residue is already
determined by its carry function, as in the proof of
Theorem~\ref{thm:main}.
Thus a bijection exists at every stage.
Any exact fixed-length code must distinguish $N_k(B)$ functions, so it
needs at least \eqref{eq:bits} bits. The dense labels attain that bound.
\end{proof}

The implementation constructs the profiles directly from the operands and
uses this ordering for ranking and unranking. The decoder selects a
representative of the equivalence class. It preserves the function rather
than the identity of the original operands.
The bit minimum is a representation result. It assigns no unit cost to
arbitrary-precision arithmetic or to evaluation of an entire truth table.

For 32-bit operands, the first exact code lengths are:
\begin{center}
\begin{tabular}{rrrr}
\toprule
Additions $k$ & Operand bits & Distinct functions & Label bits\\
\midrule
1 & 64 & 8,589,934,591 & 33\\
2 & 96 & 21,474,836,476 & 35\\
3 & 128 & 68,719,476,721 & 36\\
4 & 160 & 279,172,874,176 & 39\\
5 & 192 & 1,400,159,338,171 & 41\\
\bottomrule
\end{tabular}
\end{center}
In the third row, $N_3(2^{32})=2^{36}-15$. Thus 36 bits are sufficient
and necessary for all eight incoming-bit assignments, including the final
32-bit residue and all three carry outputs.

\section{The two-stage case and a worked finite family}
For the two-stage case~\cite{two-stage}, write $x=C+S=Bq+r$ and $r+H=Bt+v$.
The output residue is always $[v+a+b]_B$.
The possible carry functions are given in Table~\ref{tab:two}; here
$q,t\in\{0,1\}$ range independently where shown.
\begin{table}[ht]
\centering
\begin{tabular}{lll r}
\toprule
Residue & Condition & Carry outputs $(c_1,c_2)$ & Count\\
\midrule
$0\leq v\leq B-3$ & $r\leq B-2$ & $(q,t)$ & 4\\
 & $r=B-1$ & $(a,1-a)$ & 1\\
$v=B-2$ & $r\leq B-2$ & $(q,ab)$ & 2\\
 & $r=B-1$ & $(a,(1-a)\lor b)$ & 1\\
$v=B-1$ & $r\leq B-2$ & $(q,a\lor b)$ & 2\\
 & $r=B-1$ & $(a,(1-a)b)$ & 1\\
\bottomrule
\end{tabular}
\caption{Complete two-stage classification. The first residue range is empty
when $B=2$. The sum of the counts over all residues is $5B-4$.}
\label{tab:two}
\end{table}

In the first range, assign indices $2q+t$ to the constant maps and index 4
to the remaining map. Give the function label $5v$ plus this index.
In either terminal residue group, assign indices $q$ to the two ordinary
maps and index 2 to the exceptional map. The label is
$5(B-2)+3(v-B+2)$ plus this index. This gives a direct optimal codec for
every integer radix $B\geq2$.

For three stages at radix $B=8$, there are $8^4=4096$ operand tuples but
only $16\cdot8-15=113$ complete functions. There are 16 functions at each
residue $0,1,2,3,4$, and 11 at each residue $5,6,7$.
The supplied figure groups every source tuple by its exact function label.
No function is inferred from a partial set of incoming assignments.

Independently addressable local-stage functions define a different output
contract. For example, four separately recoverable byte-level two-stage
tables have $1276^4$ possibilities and need 42 bits. This product count
concerns four local tables; it is not a count of their reachable global
responses from only four initial carry pairs.

\section{Verification and related work}
An independent forward arithmetic oracle checked 54 complete radix/stage
panels: 196,160 operand tuples and 5,005,590 incoming-bit assignments.
All 13,191 distinct functions passed dense decoding, with every label
occupied exactly once. A separate two-stage sweep compared 359,996 output
triples at radices 2--24. Further checks covered 15,792 outputs at widths
8--512 bits. The tests include $k<B-1$, $k=B-1$, and $k>B-1$.
They support the implementation; the proofs establish the general statements.

Mirsalehi and Gaylord~\cite{mirsalehi} determine minimum sum-of-products
sizes for binary addition and subtraction. Their Table I contains the
expression $5\cdot2^j-4$ for a truth-table lookup statistic at bit position
$j$. Numerical agreement with our two-stage count does not identify the
objects: here the equivalence relation compares complete functions of
independent carry inputs. The present claim is the classification,
all-stage count, and dense function codec under the explicit contract
\eqref{eq:function}.

Diaconis and Fulman~\cite{diaconis} study the stochastic carry process
under random operands. Jir\'askov\'a and Okhotin~\cite{jiraskova} count
states of nondeterministic automata for positional addition of sets of
numbers. Those output contracts differ from the finite function set here.
The coefficient sequence $A_k$ itself is classical~\cite{oeis}; the proof
above establishes its role in the residue-support classification.

\begingroup
\small
\begin{thebibliography}{9}
\setlength{\itemsep}{4pt}
\bibitem{oeis}
The OEIS Foundation Inc., \textit{The On-Line Encyclopedia of Integer
Sequences}, entry A000522, ordered tuples of distinct elements.
\url{https://oeis.org/A000522}. Accessed 27 October 2026.
\bibitem{mirsalehi}
M.~Mirsalehi and T.~K.~Gaylord,
``Analytic expressions for the sizes of logically minimized truth tables
for binary addition and subtraction,''
\textit{Applied Optics} \textbf{29}(23) (1990), 3339--3344.
\href{https://doi.org/10.1364/AO.29.003339}{doi:10.1364/AO.29.003339}.
\bibitem{diaconis}
P.~Diaconis and J.~Fulman,
``Carries, shuffling, and an amazing matrix,''
\textit{American Mathematical Monthly} \textbf{116}(9) (2009), 788--803.
\href{https://doi.org/10.4169/000298909X474864}{doi:10.4169/000298909X474864}.
\bibitem{jiraskova}
G.~Jir\'askov\'a and A.~Okhotin,
``Nondeterministic state complexity of positional addition,''
\textit{Electronic Proceedings in Theoretical Computer Science}
\textbf{3} (2009), 151--161.
\href{https://doi.org/10.4204/EPTCS.3.14}{doi:10.4204/EPTCS.3.14}.
\bibitem{two-stage}
lmtlssss{\authorsymbol 🜐}, \textit{Exact State Complexity of Two-Stage Addition},
version 1.0.0, 27 October 2026.
\href{https://doi.org/10.5281/zenodo.23698617}{doi:10.5281/zenodo.23698617}.
\end{thebibliography}
\endgroup
\end{document}
